\documentclass[11pt,a4paper]{article}
\usepackage[T1]{fontenc}
\usepackage[utf8]{inputenc}
\usepackage[french]{babel}
\usepackage[a4paper,margin=2.2cm]{geometry}
\usepackage{hyperref}
\usepackage{xcolor}
\usepackage{enumitem}
\usepackage{titlesec}
\usepackage{parskip}
\usepackage{fancyhdr}

\definecolor{hmhblue}{HTML}{4A90E2}
\definecolor{hmhteal}{HTML}{50E3C2}
\definecolor{hmhdark}{HTML}{2C3E50}
\definecolor{hmhmuted}{HTML}{6B7891}

\hypersetup{colorlinks=true, linkcolor=hmhblue, urlcolor=hmhblue, citecolor=hmhblue}

\titleformat{\section}{\Large\bfseries\color{hmhdark}}{\thesection.}{0.6em}{}
\titleformat{\subsection}{\large\bfseries\color{hmhblue}}{\thesubsection}{0.6em}{}

\pagestyle{fancy}
\fancyhf{}
\fancyhead[L]{\small\textit{HearMyHands -- récap intégration web et déploiement}}
\fancyhead[R]{\small\thepage}
\renewcommand{\headrulewidth}{0pt}

\title{\vspace{-1.5cm}\textbf{HearMyHands}\\[0.2em]
       \large\color{hmhblue}{Récap de l'intégration web, des optimisations et du déploiement}}
\author{\small\textit{Pôle Interface -- Kaelig Leray}}
\date{\small Mai 2026}

\begin{document}
\maketitle
\vspace{-1.5em}

\noindent\textit{Ce document résume le travail mené sur la partie web et le déploiement
du démonstrateur HearMyHands. Il vient en complément du rapport projet principal,
qui détaille la conception scientifique des modèles. Le tout est versionné publiquement
sur \href{https://github.com/nmqx/hearmyhands}{github.com/nmqx/hearmyhands} et accessible
en ligne sur \href{https://hearmyhands.asia}{hearmyhands.asia}.}

\section{Architecture finale}

\noindent L'application est servie par un seul processus Python (Flask + Socket.IO sous
gunicorn + gevent-websocket), qui héberge à la fois l'interface utilisateur et le moteur
d'inférence. Cette fusion a été décidée après mesure: l'aller-retour HTTP entre le
serveur web et un service d'inférence séparé ajoutait $\sim$40\,ms par frame, gaspillés
sur boucle locale. Le mode service séparé reste disponible (\texttt{USE\_HTTP\_MODEL=1})
pour un déploiement distribué éventuel.

\begin{verbatim}
Navigateur --binary JPEG via WebSocket--> Flask/Socket.IO --in-process-->
   InferenceEngine (PyTorch + MediaPipe + MLP + GRU Ocarina) --ack--> Navigateur
\end{verbatim}

La logique d'inférence est isolée dans le module \texttt{HmH/inference.py} (source de
vérité unique), exposée à la fois en in-process et par un wrapper HTTP minimal
\texttt{HmH/api.py}. Le code reste 100\,\% open source.

\section{Modèles intégrés au runtime}

\subsection{Pôle Vision (Modèle 1)}
Le pôle Vision (responsable: Titouan) a fourni une partie custom basée sur
CSPDarknet53 + heatmaps + Soft-Argmax, prédisant 9 points clés du haut du corps en
entrée $416\times416$. MediaPipe Hands est appelé en complément sur des crops centrés
sur les poignets pour extraire les 21 landmarks par main. Les deux briques tournent en
parallèle (PyTorch sur GPU/CPU, MediaPipe sur CPU via un \texttt{ThreadPoolExecutor})
pour réduire la latence par frame.

\subsection{Pôle Traduction (Modèle 2)}
Deux modèles cohabitent dans le runtime, sélectionnables via un toggle UI:
\begin{itemize}[leftmargin=*,itemsep=2pt]
  \item \textbf{Mode statique}: MLP $42 \rightarrow 30 \rightarrow 23$ chargé depuis
        \texttt{Poids/\{W1,W2,b1,b2\}.json}, classification par frame des lettres
        statiques. Seuil de confiance 0,6, validation par 10 frames stables avant ajout
        au mot.
  \item \textbf{Mode dynamique}: GRU Ocarina (2 couches, hidden 128, embedding
        $42\rightarrow 256 \rightarrow 128$), entraîné par le pôle Traduction
        (responsable: Killian) sur des séquences de 45 frames. Wrist-relative + scale
        canonique $640\times480$ alignés sur \texttt{Train.py}. Reconnaît les signes
        dynamiques (J, Z) que le MLP ne peut pas traiter.
\end{itemize}

\section{Optimisations performances}

Profilage initial: 5,1\,req/s en service séparé sur \texttt{c3-highcpu-8} CPU,
avec un goulet à 196\,ms/req. Les leviers principaux:

\begin{itemize}[leftmargin=*,itemsep=2pt]
  \item Transport WebSocket en \textbf{JPEG binaire} (au lieu de base64 dataURL),
        downscale client à 480\,px, qualité 0,7 -- frames $\sim$12\,KB au lieu de 80\,KB.
  \item Pipeline non-bloquant côté client: \texttt{MAX\_IN\_FLIGHT=4} avec numéros de
        séquence, applique uniquement la dernière prédiction.
  \item Côté serveur: décodage JPEG unique (cv2 / libjpeg-turbo, plus PIL),
        pré-traitement direct numpy $\rightarrow$ tenseur GPU, \texttt{torch.inference\_mode},
        et parallélisation MediaPipe $\leftrightarrow$ PyTorch.
  \item Fusion in-process Flask + InferenceEngine pour éliminer l'overhead HTTP local.
\end{itemize}

\noindent\textbf{Résultat:} 44\,req/s sustained en in-process sur la même machine
($\sim$23\,ms/req), soit $\sim8,\!5\times$ plus rapide. Le client est ensuite cappé à
30\,fps (au-delà, le serveur CPU ne suit pas et les frames seraient droppées).

\section{Mode apprentissage (style Anki)}

Sous-section \texttt{/learn} ajoutée pour la dimension pédagogique du projet
(\textbf{ODD~4}). Une fiche par lettre montre la vidéo modèle (signée par Nestor pour la
majorité, Louna pour le Q) à côté de la webcam de l'apprenant. La pipeline d'inférence
existante est réutilisée: dès que la prédiction MLP correspond à la lettre cible
pendant 5 frames stables au-dessus de 60\,\% de confiance, une validation visuelle
``Bien joué!'' s'affiche. L'apprenant note ensuite sa carte parmi
\textit{à revoir} / \textit{j'apprends} / \textit{je connais}, et un algorithme de
répétition espacée gère la file (intervalles 3, 5, 10, 20, 40, 80 cartes). L'état est
persistant via \texttt{localStorage}, et une bibliothèque expose toutes les lettres et
permet de remettre en révision celles déjà connues.

\section{Déploiement}

Le démonstrateur est hébergé sur Google Cloud Compute (\texttt{c3-highcpu-8} Debian 12
en \texttt{europe-west1-b}, IP statique réservée), choisi pour la proximité réseau avec
les utilisateurs français ($\sim$24\,ms de RTT depuis Nantes, vs $\sim$120\,ms quand le
serveur était hébergé en \texttt{us-central1}). nginx assure le reverse proxy avec les
headers WebSocket adéquats, gunicorn fait tourner l'app avec un worker
\texttt{geventwebsocket} pour Socket.IO. TLS via Let's Encrypt + certbot
(auto-renouvellement par timer systemd), avec une seule chaîne couvrant les quatre
\textit{server names}: \texttt{hearmyhands.asia}, \texttt{www.hearmyhands.asia},
\texttt{nmqx.lol}, \texttt{www.nmqx.lol}. Cloudflare est posé en frontal pour la
protection anti-bots, le cache statique et le HTTP/3.

Le service \texttt{hmh.service} est piloté par systemd (\texttt{Restart=on-failure})
et un workflow de redéploiement tient en une commande:
\texttt{ssh hmh@vm "cd hearmyhands \&\& git pull \&\& sudo systemctl restart hmh"}.

\section{Monitoring et outillage développeur}

Un dashboard temps-réel \texttt{/monitor} polle \texttt{/stats} chaque seconde et affiche
CPU total + par cœur, RAM, disque, load average, ainsi que la consommation du worker
gunicorn. Implémentation: psutil + un template HTML autonome, dans la même DA glass +
gradient. Le mode sombre (toggle lune/soleil dans la navbar, persistant via
\texttt{localStorage}) couvre l'ensemble du site, par défaut le mode clair.

Le module \texttt{HmH/sign\_classifier.py} prévoit un chargement gracieux des poids
Ocarina: en l'absence du fichier \texttt{ocarina\_gru\_v1.pth}, l'endpoint
\texttt{/sign\_predict} retourne 503 et la webapp désactive silencieusement les
prédictions dynamiques. Cela autorise des démos même lorsque la branche temporelle
n'est pas embarquée.

\section{Liens utiles}

\begin{itemize}[leftmargin=*,itemsep=2pt]
  \item Application live: \href{https://hearmyhands.asia/translate}{hearmyhands.asia/translate}
  \item Apprentissage: \href{https://hearmyhands.asia/learn}{hearmyhands.asia/learn}
  \item Monitoring: \href{https://hearmyhands.asia/monitor}{hearmyhands.asia/monitor}
  \item Code source: \href{https://github.com/nmqx/hearmyhands}{github.com/nmqx/hearmyhands}
  \item ODD 10: \href{https://www.un.org/sustainabledevelopment/fr/inequality/}{un.org/sustainabledevelopment/fr/inequality/}
  \item ODD 4: \href{https://www.un.org/sustainabledevelopment/fr/education/}{un.org/sustainabledevelopment/fr/education/}
  \item Laboratoire LS2N: \href{https://www.ls2n.fr/}{ls2n.fr}
\end{itemize}

\vspace{0.6em}
\noindent\textit{Remerciements: à Hélène Pérennou pour le suivi pédagogique, à Matthieu
Perreira Da Silva pour l'accompagnement technique côté LS2N, et aux 11 autres membres
de l'équipe pour le travail collectif.}

\end{document}
